% =================================================================
% 第6章 总结与展望
% =================================================================
\chapter{总结与展望}
\label{chap:conclusion_and_future_work}

\section{全文工作总结}
\label{sec:conclusion_summary}

随着工业 4.0 与智能制造的深入推进，构建高保真、高实时性的物理实体虚拟副本已成为工业数字孪生技术落地的核心诉求。
近年来，以三维高斯溅射（3D Gaussian Splatting, 3DGS）为代表的显式辐射场技术，凭借其卓越的渲染帧率与照片级画质，为新视角合成带来了范式级的突破。
然而，当直接将其应用于包含大量高反光金属、纯色弱纹理面板以及复杂设备遮挡的类工业数字孪生场景时，传统 3DGS 暴露出严重的“几何坍塌”、“浮空伪影”以及“显存消耗巨大”等致命缺陷，难以满足端侧设备的大规模部署需求。
针对上述挑战，本文以实现面向工业孪生场景的高效、高保真虚实映射为目标，围绕 3DGS 的空间几何约束增强与模型轻量化压缩展开了系统性研究，并最终构建了端到端的融合算法原型验证系统。

本文的主要研究工作与创新成果总结如下：

\textbf{（1）提出了面向类工业场景的多模态自适应几何约束重建算法}

针对工业设备表面强烈的各向异性反射与弱纹理区域导致的光度一致性失效问题，本文打破了纯视觉优化的单一路径局限，提出了一种融合多模态几何先验的自适应重建算法。
首先，通过硬件级时空对齐，将 LiDAR 绝对尺度点云与视觉大模型（Depth Anything）推断的稠密相对深度场进行跨模态配准，为 3DGS 注入了极其稳健的初始物理骨架。
其次，设计了基于深度局部高频特征感知的自适应置信度掩膜（Adaptive Mask）机制，动态识别场景中的平坦区与高频边缘。
最后，构建了联合 Huber 深度惩罚与一阶法向平滑正则化的联合损失函数，在掩膜的动态调度下，强制高斯基元贴合真实的物理表面。
实验表明，该算法彻底清除了金属表面的絮状浮空噪点，完美修复了纯色机柜的几何撕裂，在大幅降低深度均方根误差与倒角距离（Chamfer Distance）的同时，完好保留了机加工纹理的高频细节。

\textbf{（2）提出了面向典型孪生场景的 3DGS 轻量化表达与端侧渲染优化算法}

针对 3DGS 庞大的高维参数体量导致端侧设备显存溢出（OOM）与渲染卡顿的工程痛点，本文从数据结构重构与渲染管线优化双管齐下。
在几何层面，构建了融合透射率、全局视图可见性以及空间尺度的多维度重要性评估指标，结合长尾分布自适应阈值，精准剥离了隐藏在不可见区域的海量冗余“幽灵高斯”。
在属性层面，引入矢量量化（Vector Quantization, VQ）技术，将占据超过 80\% 存储开销的高维球谐函数（SH）系数映射为紧凑的离散码本，并结合直通估计器（STE）完成量化感知微调（QAT），有效避免了硬量化带来的色彩断层。
此外，本文深度重构了 Web 端的渲染管线，利用 WebGPU 的计算着色器实现显存内硬件级解码，并适配移动端 TBDR 架构引入了局部排序与光线提前终止机制，彻底打通了轻量化模型在普通轻薄本与移动平板上的实时流畅漫游路径。

\textbf{（3）构建了融合改进 3DGS 的算法管线并完成系统级综合验证}

本文将上述“多模态几何约束”与“轻量化剪枝量化”在底层逻辑与训练管线中进行了深度融合，提出了一套分阶段的渐进式训练策略（几何预热、动态生长与剪枝、色彩微调）。
研究深入揭示了两者之间的“协同效应”：稳固的物理几何底座不仅提升了渲染保真度，更为大刀阔斧的剪枝手术提供了精准的安全向导，避免了传统盲剪导致的画面破损。
在自主构建的复杂机房孪生数据集上的综合测试表明，相较于原始 3DGS，本文融合算法在保持 31.20 dB 极高 PSNR 水准的前提下，将场景模型的物理存储体积极限压缩了约 95.8\%（降至 48.2 MB），峰值渲染显存断崖式降至 0.45 GB，并在常规硬件环境下实现了 120 FPS 的极致渲染帧率。
这一成果无可辩驳地验证了本文方法在破除算力壁垒、赋能实际工业数字孪生业务系统方面的巨大潜力与工程落地价值。

\section{研究不足与未来展望}
\label{sec:future_work}

尽管本文在面向类工业孪生场景的 3DGS 高保真重建与轻量化方面取得了一定的成果，
但受限于目前的技术瓶颈以及实验硬件的算力制约，
本研究仍存在若干局限性。
在工业元宇宙与下一代数字孪生技术不断演进的宏观背景下，
未来的研究工作可以从以下四个维度进行更深入的探索：

\textbf{（1）面向动态工业场景的 4D 高保真重建}

本文目前的算法管线主要针对静态的机房与设备环境进行三维重建，
其底层假设是场景在数据采集过程中保持绝对静止。
然而，
真实的工业车间往往存在大量高度动态的元素，
例如运转的机械臂，
移动的 AGV 小车，
以及穿梭的工作人员。
未来可引入时间维度，
探索 4D 高斯溅射（4D Gaussian Splatting）技术，
结合非刚性形变场（Deformation Field），
实现工业设备动态运行过程的实时高保真孪生。

\textbf{（2）基于物理的材质解耦与实时重光照}

本课题利用球谐函数（SH），
将复杂的全局光照和材质反射属性高度耦合地“烘焙”在了离散的高斯基元中。
一旦数字孪生环境中的光源发生变化，
当前模型无法做出正确的动态响应。
未来研究可致力于 3DGS 的本征分解（Intrinsic Decomposition），
结合基于物理的渲染（Physically Based Rendering），
将高斯球的颜色拆分为基础反照率（Albedo），
粗糙度（Roughness），
和金属度（Metallic），
从而支持在虚拟孪生空间中进行任意光源位置和强度的“重光照”编辑。

\textbf{（3）语义级孪生与物联网数据的深度融合}

目前的改进算法主要集中在“视觉”与“几何”层面的高保真还原，
模型本质上仍是一个缺乏语义理解的渲染黑盒。
工业数字孪生的最终业务目的是交互与精准监控。
未来可以探索将二维基础视觉大模型与 3DGS 进行跨维度的特征蒸馏，
实现三维空间中高斯球的“实例级语义分割”。
在此基础上，
将现实世界中的物联网（IoT）传感器实时数据精准锚定，
并挂载到对应的三维语义高斯模型上，
实现真正意义上的业务级虚实映射。

\textbf{（4）超大规模厂区的分布式协同重建架构}

受限于单机单卡的物理显存与算力，
本文目前的轻量化与训练管线仅在单一机房级别的中等尺度场景上进行了实验验证。
对于面积达数万平方米的超大型重工业园区，
端到端的单次训练依然存在不可逾越的瓶颈。
未来可研究基于空间区块划分的分布式协同训练策略，
探索更加高效的场景分块合并与边界光度过渡算法，
以支撑无界大规模工厂的全要素数字孪生底座构建。